%==============================================================================
%  RELATÓRIO 1
%  Entrega: 24/03/2026
%==============================================================================

\section{Introdução}

O presente relatório constitui a primeira etapa do projeto prático da disciplina
PRO3341 -- Modelagem e Otimização de Sistemas de Produção. O objetivo do projeto
é aplicar técnicas de Pesquisa Operacional a um problema real identificado em
parceria com uma empresa, produzindo resultados quantitativamente embasados e
comparáveis à prática vigente.

A empresa escolhida para o projeto é a Rentex Indústria Têxtil LTDA, fabricante
de tecidos de malha e planos localizada no interior do estado de São Paulo. O
problema proposto é a otimização do mix de produção semanal: determinar quais
produtos fabricar e em que quantidade, respeitando as capacidades instaladas dos
recursos produtivos, de modo a maximizar a margem de contribuição total.

Este relatório apresenta a caracterização da empresa e de seu ambiente produtivo,
a definição e justificativa do problema, uma concepção preliminar do modelo
matemático e a revisão bibliográfica de suporte, que inclui um artigo indexado
diretamente relacionado ao tema e a resolução de um exercício do livro-texto
com auxílio de software de otimização.

%==============================================================================
\section{Caracterização da Empresa}
%==============================================================================

\subsection{Histórico e Atuação}

A Rentex Indústria Têxtil LTDA é uma empresa de médio porte fundada na década de
1990, com sede em Americana, SP, município historicamente reconhecido como polo
têxtil brasileiro \cite{abit2023}. A empresa atua no segmento de tecidos para
confecção, produzindo malhas e tecidos planos destinados a fabricantes de
vestuário de médio e alto padrão nas regiões Sudeste e Sul do país.

O contato institucional para o projeto é Pedro Charmillot Silva, Diretor de
Novos Negócios, cuja carta de aprovação encontra-se em anexo a este relatório.

\subsection{Portfólio de Produtos}

O portfólio ativo da Rentex é composto por seis famílias de tecidos, cada uma
com características distintas de processo, custo e posicionamento de mercado:

\begin{enumerate}[label=\textbf{(\arabic*)}]
  \item \textbf{Malha Jersey} -- tecido de malha leve, utilizado em camisetas e
        peças básicas de vestuário. É o produto de maior volume da empresa, com
        demanda estável ao longo do ano.

  \item \textbf{Malha Rib} -- malha elástica com estrutura de canelado, empregada
        em golas, punhos e peças esportivas. Exige ajuste mais preciso nos teares
        circulares.

  \item \textbf{Tecido Plano Popeline} -- tecido leve de algodão ou misto,
        utilizado em camisas sociais e blusas. Requer processamento nos teares de
        rapier e etapa completa de engomagem.

  \item \textbf{Tecido Plano Oxford} -- variação mais encorpada do popeline, com
        trama característica. Produto de maior valor agregado dentro dos planos,
        voltado ao segmento de moda feminina.

  \item \textbf{Denim Leve} -- tecido jeans de gramatura reduzida (abaixo de
        $12$ oz/m²), produzido para calças e jaquetas de linha premium. Consome a
        maior quantidade de fio por unidade produzida.

  \item \textbf{Tecido Técnico} -- tecido funcional com tratamento de acabamento
        especial (repelência, antichama ou UPF), destinado ao segmento de
        uniformes corporativos e equipamentos de proteção individual (EPI).
        Apresenta a maior margem unitária do portfólio.
\end{enumerate}

\subsection{Estrutura Produtiva}

O parque fabril da Rentex opera em regime de dois turnos (segunda a sábado) e
conta com os seguintes recursos principais:

\begin{itemize}
  \item \textbf{Teares:} 20 teares circulares (para malhas) e 8 teares de rapier
        (para planos e denim), com disponibilidade semanal total de
        aproximadamente 1.200 horas-máquina.

  \item \textbf{Tinturaria:} setor de tingimento por esgotamento e foulard, com
        capacidade de 800 horas-equipamento por semana. É o recurso que concentra
        os maiores custos variáveis do processo.

  \item \textbf{Acabamento:} linha de rameuse, compactadora e inspeção final, com
        600 horas disponíveis por semana. Constitui o gargalo mais frequente para
        os tecidos técnicos, cujo tratamento de superfície é mais demorado.

  \item \textbf{Matéria-prima (fio):} estoque regulador semanal de 15.000 kg de
        fio, composto por algodão cardado, poliéster e misturas. A reposição é
        feita semanalmente por fornecedores fixos, tornando esse limite efetivo no
        horizonte de planejamento.

  \item \textbf{Mão de obra direta (MOD):} 900 horas de mão de obra direta
        disponíveis por semana, distribuídas entre os setores.
\end{itemize}

\subsection{Contexto Competitivo}

O setor têxtil brasileiro enfrenta pressão crescente de importados asiáticos,
especialmente na categoria de malhas básicas. A resposta estratégica das empresas
nacionais de médio porte tem sido a especialização em produtos de maior valor
agregado e a redução de ineficiências operacionais
\cite{iemi2022}.

Nesse contexto, a decisão sobre o mix de produção semanal tem impacto direto na
rentabilidade: produtos com maior margem de contribuição consomem mais recursos
escassos, e a alocação equivocada entre famílias de produto pode representar
perda de dezenas de milhares de reais por semana.

\newpage
